\documentclass[12pt,a4paper]{article}
\usepackage[utf8]{inputenc}
\usepackage[T5]{fontenc}
\usepackage[vietnamese]{babel}
\usepackage{graphicx}
\usepackage{amsmath}
\usepackage{hyperref}
\usepackage{geometry}
\geometry{a4paper, margin=1in}

\title{\textbf{Hướng dẫn Sử dụng Chi tiết: Viettel Auto-Tuner\\(Hệ thống Tối ưu hóa tham số bộ điều khiển trực thăng 2-DOF)}}
\author{}
\date{}

\begin{document}
\maketitle

\section{Giới thiệu Tổng quan}
Tài liệu này cung cấp hướng dẫn vận hành toàn diện \textbf{Viettel Auto-Tuner}, từ quá trình cài đặt, khởi chạy hệ thống, cho đến các thao tác cấu hình, giám sát và xử lý lỗi chuyên sâu.

\section{Hướng dẫn Khởi động Hệ thống (Startup \& Installation)}

\textbf{Điều kiện cần:}
\begin{itemize}
    \item Cài đặt Python 3.8+ và Node.js 18+ trên máy tính.
    \item Khởi tạo sẵn môi trường ảo Python tại thư mục \texttt{backend/venv}.
\end{itemize}

\textbf{Các bước khởi động tự động (Khuyến nghị):}
\begin{enumerate}
    \item Truy cập thư mục gốc của dự án (ví dụ: \texttt{d:\textbackslash UIT\textbackslash Research\textbackslash Viettel\textbackslash}).
    \item Khởi chạy tệp thực thi \texttt{start\_project.bat}.
    \item Kiểm tra 2 cửa sổ Command Prompt riêng biệt cho Backend (cổng \texttt{8088}) và Frontend (cổng \texttt{3088}). Việc khởi chạy thành công được minh họa như trong \textbf{Hình~\ref{fig:startup_terminals}}.
    \item Truy cập giao diện web tại địa chỉ \texttt{http://localhost:3088} thông qua trình duyệt.
\end{enumerate}

\textbf{Các bước khởi động thủ công:}
\begin{enumerate}
    \item Mở Command Prompt cho Backend, chạy tuần tự các lệnh sau:
    \begin{itemize}
        \item \texttt{cd backend}
        \item \texttt{venv\textbackslash Scripts\textbackslash activate}
        \item \texttt{python -m uvicorn main:app --port 8088}
    \end{itemize}
    \item Mở Command Prompt thứ hai cho Frontend, chạy các lệnh sau:
    \begin{itemize}
        \item \texttt{cd frontend}
        \item \texttt{npm run dev -- --port 3088}
    \end{itemize}
    \item Truy cập giao diện web tại địa chỉ \texttt{http://localhost:3088} thông qua trình duyệt.
\end{enumerate}

\begin{figure}[htbp]
    \centering
    \includegraphics[width=0.85\textwidth]{images/startup_terminals.png}
    \caption{Cửa sổ Command Prompt khi khởi động Backend và Frontend thành công}
    \label{fig:startup_terminals}
\end{figure}

\textbf{Kết quả kỳ vọng:} \\
Hệ thống Frontend hiển thị giao diện chính của Viettel Auto-Tuner. Cửa sổ Backend báo hiệu "Application startup complete".

\section{Phân tích Tối ưu hóa (Ideal Tuning)}

\textbf{Điều kiện cần:}
\begin{itemize}
    \item Hoàn tất bước khởi động hệ thống.
    \item Đảm bảo Backend không báo lỗi kết nối.
\end{itemize}

\textbf{Các bước thực hiện:}
\begin{enumerate}
    \item Mở thẻ \textbf{[Experiments]}, chọn thẻ phụ \textbf{[1. Ideal Tuning]}.
    \item Thiết lập điểm đặt mục tiêu tại bảng \textbf{[Global Config]} (ví dụ: Setpoint Pitch, Setpoint Yaw).
    \item Lựa chọn hàm mục tiêu tại danh sách \textbf{[Objective Function Structure]}.
    \item Chọn chiến lược phạt tại danh sách \textbf{[Tuning Profile (Weights)]} (nhập \texttt{1. Balanced}, \texttt{2. Aggressive Tracking}, \texttt{3. Eco \& Smooth}, hoặc \texttt{4. Strict Safety}). Tham khảo \textbf{Hình~\ref{fig:ideal_tuning_config}}.
    
\begin{figure}[htbp]
    \centering
    \includegraphics[width=0.85\textwidth]{images/ideal_tuning_config.png}
    \caption{Thẻ [1. Ideal Tuning] với các thông số cấu hình Global Config và Tuning Profile}
    \label{fig:ideal_tuning_config}
\end{figure}

    \item Nhập số lượng thế hệ tối ưu tại các trường như \textbf{[GA Iters]}, \textbf{[PSO Iters]} (ví dụ: \texttt{50}).
    \item Nhấp chuột vào nút \textbf{[Run Analysis]}. Quan sát thanh tiến trình báo cáo phần trăm hoàn thành thuật toán (\textbf{Hình~\ref{fig:ideal_tuning_run}}).

\begin{figure}[htbp]
    \centering
    \includegraphics[width=0.85\textwidth]{images/ideal_tuning_run.png}
    \caption{Nút [Run Analysis] và thanh tiến trình đang thực thi}
    \label{fig:ideal_tuning_run}
\end{figure}

    \item Nhấp chuột vào các nút thuật toán trên biểu đồ \textbf{Main Comparison} để bật hoặc tắt hiển thị. So sánh dữ liệu hội tụ tại \textbf{Hình~\ref{fig:ideal_tuning_charts}}.

\begin{figure}[htbp]
    \centering
    \includegraphics[width=0.85\textwidth]{images/ideal_tuning_charts.png}
    \caption{Đồ thị ``Comparison of Best Models'' và ``Convergence Chart''}
    \label{fig:ideal_tuning_charts}
\end{figure}

\end{enumerate}

\textbf{Kết quả kỳ vọng:} \\
Biểu đồ thời gian (Time Response) cập nhật đường tín hiệu của các thuật toán. Các bảng thông số hiệu năng (Rise Time, Overshoot) hiển thị giá trị cụ thể.

\section{Đánh giá Độ bền vững (Robustness Sweep)}

\textbf{Điều kiện cần:}
\begin{itemize}
    \item Hoàn thành "Phân tích Tối ưu hóa" ít nhất một thuật toán để có tham số tham chiếu.
\end{itemize}

\textbf{Các bước thực hiện:}
\begin{enumerate}
    \item Mở thẻ \textbf{[Experiments]}, chọn thẻ phụ \textbf{[2. Robustness Sweep]}.
    \item Nhập số lượng bài kiểm tra tại trường \textbf{[Number of Configs]} (ví dụ: \texttt{6}).
    \item Điền giá trị nhiễu ban đầu vào cột \textbf{[Base Value]} (ví dụ: \texttt{Wind Pitch = 0.01}).
    \item Điền bước nhảy nhiễu vào cột \textbf{[Step Increment]} (ví dụ: \texttt{0.01} mỗi bước). Tham khảo \textbf{Hình~\ref{fig:robustness_config}}.

\begin{figure}[htbp]
    \centering
    \includegraphics[width=0.85\textwidth]{images/robustness_config.png}
    \caption{Bảng cấu hình ``Robustness Analysis'' và các giá trị Base Value, Step Increment}
    \label{fig:robustness_config}
\end{figure}

    \item Nhấp chuột vào nút \textbf{[Run Sweep]}. 
    \item Đánh giá kết quả trên các biểu đồ (Rise Time, Overshoot, Error) theo từng mức nhiễu (\textbf{Hình~\ref{fig:robustness_charts}}).

\begin{figure}[htbp]
    \centering
    \includegraphics[width=0.85\textwidth]{images/robustness_charts.png}
    \caption{Đồ thị phân tích độ bền vững (Rise Time, Overshoot, Error)}
    \label{fig:robustness_charts}
\end{figure}

\end{enumerate}

\textbf{Kết quả kỳ vọng:} \\
Hệ thống kết xuất 5 biểu đồ đường (Line chart) thể hiện sự suy giảm hiệu năng của các bộ điều khiển theo hàm tăng dần của nhiễu động.

\section{Giám sát trên Mô phỏng 3D}

\textbf{Điều kiện cần:}
\begin{itemize}
    \item Ít nhất một thuật toán đã được tối ưu hóa thành công thông qua tab Ideal Tuning.
\end{itemize}

\textbf{Các bước thực hiện:}
\begin{enumerate}
    \item Mở thẻ \textbf{[3D Simulation]} trên thanh điều hướng chính.
    \item Lựa chọn thuật toán tại danh sách thả xuống \textbf{[Controller Algorithm]} (ví dụ: \texttt{PSO}). Tham khảo \textbf{Hình~\ref{fig:3d_sim_settings}}.

\begin{figure}[htbp]
    \centering
    \includegraphics[width=0.85\textwidth]{images/3d_sim_settings.png}
    \caption{Bảng cấu hình ``3D SIMULATION SETTINGS'' và việc chọn Controller Algorithm}
    \label{fig:3d_sim_settings}
\end{figure}

    \item Kéo thả khối \textbf{[Joystick]} ở góc dưới bên phải màn hình để thay đổi điểm đặt của trực thăng.
    \item Giám sát thông số góc độ tại bảng \textbf{TELEMETRY} ở góc trên bên phải (\textbf{Hình~\ref{fig:3d_sim_dashboard}}).

\begin{figure}[htbp]
    \centering
    \includegraphics[width=0.85\textwidth]{images/3d_sim_dashboard.png}
    \caption{Giao diện mô phỏng 3D với bảng Telemetry và Joystick}
    \label{fig:3d_sim_dashboard}
\end{figure}

\end{enumerate}

\textbf{Kết quả kỳ vọng:} \\
Mô hình trực thăng 3D di chuyển khớp với điểm đặt Joystick theo thời gian thực. Bảng Telemetry liên tục cập nhật dữ liệu với độ trễ thấp.

\section{Xử lý Sự cố}

\subsection{Cổng giao tiếp bị chiếm dụng (Port in use)}
\begin{itemize}
    \item \textbf{Nguyên nhân:} Cổng \texttt{8088} hoặc \texttt{3088} đang bị một ứng dụng ngầm sử dụng.
    \item \textbf{Giải pháp:} 
    \begin{enumerate}
        \item Mở Task Manager.
        \item Kết thúc (End task) toàn bộ các tiến trình \texttt{python.exe} và \texttt{node.exe}.
        \item Khởi động lại \texttt{start\_project.bat}.
    \end{enumerate}
\end{itemize}

\subsection{Bão hòa điện áp (Voltage Saturation)}
\begin{itemize}
    \item \textbf{Nguyên nhân:} Tín hiệu ngõ ra vượt ngưỡng vật lý $\pm 24\text{V}$ do thuật toán yêu cầu tốc độ phản hồi quá cao.
    \item \textbf{Giải pháp:} 
    \begin{enumerate}
        \item Mở thẻ \textbf{[1. Ideal Tuning]}.
        \item Chọn lại mục \textbf{[Tuning Profile (Weights)]} sang \textbf{[4. Strict Safety]}.
        \item Nhấp chuột vào nút \textbf{[Run Analysis]} để tối ưu lại.
    \end{enumerate}
\end{itemize}

\subsection{Thuật toán không hội tụ}
\begin{itemize}
    \item \textbf{Nguyên nhân:} Các thuật toán tối ưu (GA, PSO) bị kẹt tại cực tiểu cục bộ do cấu hình Iters quá nhỏ.
    \item \textbf{Giải pháp:} 
    \begin{enumerate}
        \item Mở thẻ \textbf{[1. Ideal Tuning]}.
        \item Thay đổi giá trị tại trường \textbf{[Iters]} lên mức tối thiểu \texttt{50}.
        \item Nhấp chuột vào nút \textbf{[Run Analysis]}.
    \end{enumerate}
    Lưu ý: Nếu giao diện 3D báo lỗi đỏ ``Controller has not been tuned yet'' (\textbf{Hình~\ref{fig:troubleshoot_untuned}}), bắt buộc phải chạy \textbf{[Run Analysis]} thành công trước.
\end{itemize}

\begin{figure}[htbp]
    \centering
    \includegraphics[width=0.85\textwidth]{images/troubleshoot_untuned.png}
    \caption{Cảnh báo đỏ do Controller chưa được tune}
    \label{fig:troubleshoot_untuned}
\end{figure}

\subsection{Giao diện bị sập (Crash trình duyệt)}
\begin{itemize}
    \item \textbf{Nguyên nhân:} Biểu đồ ECharts render quá nhiều lịch sử vòng lặp (hàng ngàn điểm dữ liệu) gây tràn bộ nhớ DOM.
    \item \textbf{Giải pháp:} 
    \begin{enumerate}
        \item Nhấn phím \texttt{F5} để tải lại giao diện.
        \item Tắt bớt các hiển thị lịch sử không cần thiết. Không lạm dụng nút \textbf{[Show All]}.
    \end{enumerate}
\end{itemize}

\end{document}
